% Minimal paper sections (draft, 2026-11-05): M21, unlearning, subliminal learning, emergent misalignment.
% Numbers: experiments/training-run/notes/decisions.md (2026-10-07 .. 10-20) and experiments/em/notes/literature.md.
\documentclass[11pt]{article}
\usepackage{amsmath,amssymb}
\usepackage{booktabs}
\usepackage[margin=1in]{geometry}
\usepackage{hyperref}
\usepackage[numbers]{natbib}
\newcommand{\Mtwo}{M21}

\begin{document}

\section{The pairing gap (\Mtwo)}
\label{sec:m21}

Let $\{(q_i,a_i)\}_{i=1}^{n}$ be question--answer pairs and $S_{ij}=-\log p_\theta(a_j\mid q_i)$ the summed
loss of answer $a_j$ after question $q_i$ under the model $\theta$, with $D$ the number of answer tokens.
For a cyclic rotation $\pi$ of the answers,
\begin{equation}
  \Delta_\pi \;=\; \frac{1}{D}\sum_{i=1}^{n}\bigl(S_{i\pi(i)}-S_{ii}\bigr),
  \qquad
  \bar\Delta \;=\; \mathbb{E}_\pi\,\Delta_\pi ,
\end{equation}
in nats per answer token, estimated with four seeded rotations and no training. Writing
$S_{ij}=r_i+c_j+e_{ij}$, the row and column effects cancel for every $\pi$, and averaging over the $n-1$
rotations gives $\bar\Delta=-\tfrac{n}{n-1}\tfrac{1}{D}\sum_i e_{ii}$: \Mtwo{} is the interaction between
questions and their own answers. It is invariant to any change that depends only on the question (a
broken prompt format, a suppression of the topic) or only on the answer (vocabulary, length, a penalty on
forbidden outputs). In description-length terms, $D\bar\Delta$ is the code length the model already
saves on the pairings before training, the head start of the elicitation curve of
\citet{donoway2026bits} at $n=0$. By Jensen's inequality the $n$-way identification information of any
read-out is bounded by $\bar\Delta$ per item, so a gap at its floor means no read-out can pick the
answers above chance. We read \Mtwo{} against two references: a model that cannot hold the pairings sets
the floor (TinyStories-1B on WMDP-bio reads $0.22$), and on free text, where absolute values do not
transfer across tasks, we report $r=\bar\Delta_B/\bar\Delta_{\mathrm{own}}$, the fraction of the gap the
model shows on its own answers to the same questions; $r\ge 0.5$ reads elicit-side. \Mtwo{} is a
sufficient sign of held information, not a necessary one: information that no conditional likelihood
reflects reads $0$ and can still be elicited.

\section{Unlearning: the capability stays in the weights}
\label{sec:unlearning}

We read \Mtwo{} on held-out WMDP-bio facts \citep{li2024wmdp} for Zephyr-7B and six unlearned versions,
and for the Deep Ignorance suite \citep{obrien2025deepignorance}. Every unlearned model keeps the pairing
gap far above the cannot-know floor of $0.22$ nats/token: RMU $1.30$, ELM $1.24$, GradDiff $0.96$,
RMU+LAT $1.90$, NPO $2.52$ and SimNPO $2.64$, against $1.58$ for the original. NPO, whose multiple-choice
read-out sits at chance, shows one of the largest gaps: its loss on the bio answers rose by about
$244$ nats per token, a question-independent shift that \Mtwo{} removes by construction, while the
log-odds between true and other answers were left intact. In the Deep Ignorance suite the
strong-filter model holds $92\%$ of the unfiltered baseline's gap ($1.23$ vs $1.34$), and the
circuit-breaker models hold $60$--$66\%$ ($0.79$--$0.89$).

Relearning confirms the reading. Fine-tuning on $32$ true facts of a disjoint split raises \Mtwo{} on the
held-out facts above the untrained baseline for every model (e.g.\ circuit breakers $0.79\to1.47$,
circuit breakers with latent adversarial training $0.87\to1.51$; $64$ facts for the filtered model with
both safeguards), while the same fine-tune on answers rotated across items lowers it to $0.23$--$0.40$.
Separating the description-length cost into a fact-free part (the same sweep on rotated answers) and a
fact-specific part shows that the rising phase of the EDL curve in the adversarially trained safeguards
is paid with or without the facts: a gate, crossed after $128$--$573$ facts, not teaching. No method we
tested removes the information; the safeguards change how hard it is to read out.

\section{Subliminal learning is elicitation}
\label{sec:sl}

In subliminal learning (SL) a student trained on trait-free data from a teacher acquires the teacher's
trait \citep{cloud2025subliminal}. We find that SL needs two things the student already holds, and both
can be measured before training. First, the data must read elicit-side: on number sequences from
Qwen2.5 teachers with an animal-preference prompt, SL appears only when teacher and student are the same
model (target share $+0.73$ to $+0.98$, and $+0.26$/$+0.85$ for leopard), weakly under a shared
initialisation ($-0.03$ to $+0.10$) and not across models ($\le 0.02$). \Mtwo{} of the neutral teacher's
numbers separates these pairs perfectly (same model $0.91$--$1.00$ nats/token, other models
$0.10$--$0.52$; AUC $1.00$). On the letter-counting transfer sets of \citet{dubinski2026beyondowls},
the Qwen student reads elicit-side ($r=0.63$) where SL was reported and the Nemotron student teach-side
($r=0.35$) where it was not, in $4$ of $5$ checkable cases. Second, the data must carry a trait the
student holds: the trait gap, the loss reduction the trait prompt gives the student on the data, net of
the same on neutral data, separates transfer from no transfer (all transferring cases $\ge0.15$, all
others $\le0.047$; AUC $1.00$). \Mtwo{} does not see the trait itself: trait-free control data from the
same teacher reads exactly as elicit-side and transfers nothing.

\subsection{Graded transfer across model families: Logit-Linear Selection}
\label{sec:lls}

Logit-Linear Selection (LLS) \citep{adenali2026lls} passes a trait without generating any data. A
teacher scores pairs from a public preference dataset by how much a trait system prompt $s$ shifts its
preference for the chosen answer $y^+$ over the rejected one $y^-$,
\begin{equation}
  w \;=\; \bigl[\log p(y^+\mid s,x)-\log p(y^-\mid s,x)\bigr]-\bigl[\log p(y^+\mid x)-\log p(y^-\mid x)\bigr],
\end{equation}
normalised by the pair's token count, and the highest-scoring pairs become a DPO training set. Neither
answer mentions the trait. $w$ is our trait gap computed on the teacher; we compute it on each
\emph{student}, before training, and ask whether it ranks how much of the trait the student acquires.

\paragraph{Setup.} Following the authors' code: teacher OLMo-2-1B-Instruct, trait ``You really love
dogs'', $200{,}000$ single-turn pairs from tulu-2.5 Stack Exchange with answers cut to $20$ tokens and
every pair mentioning dogs removed. $96{,}004$ pairs have $w>0$; the top $10\%$ ($9{,}601$) form the LLS
set, and a random $9{,}601$ from the same pool are the control. Each student is trained with DPO on each
set (LoRA rank $64$, learning rate $5\times10^{-4}$, $\beta=0.1$, batch $512$, $10$ passes). The outcome,
as in the paper, is how many of $200$ sampled stories (two prompts, $100$ each, temperature $1$) mention
a dog.

\begin{table}[h]
\centering
\small
\begin{tabular}{llrrrr}
\toprule
student & relation to teacher & trait gap (nats/token) & after LLS set & after random set & Fisher $p$ \\
\midrule
OLMo-2-1B-Instruct & the teacher itself & $+0.329$ & $47$ & $8$ & $<10^{-4}$ \\
OLMo-2-7B-Instruct & same family & $+0.122$ & $6$ & $0$ & $0.03$ \\
Qwen2.5-7B-Instruct & different family & $+0.073$ & $9$ & $0$ & $0.004$ \\
Qwen3-8B & different family & $+0.055$ & $2$ & $1$ & $1.0$ \\
Olmo-3-7B-Instruct & newer generation & $+0.017$ & $7$ & $5$ & $0.77$ \\
rnj-1-instruct & different family & $+0.003$ & $2$ & $4$ & $0.69$ \\
\bottomrule
\end{tabular}
\caption{Dog preference by Logit-Linear Selection. Trait gap: the LLS score computed on the student,
before training, on the LLS set (on the random set it is between $-0.007$ and $+0.003$ for every
student). Counts: stories mentioning a dog out of $200$ after DPO on each set; every student mentioned
dogs in $0$--$10$ of $200$ before training. Fisher: one-sided, LLS vs random.}
\label{tab:lls}
\end{table}

\paragraph{Result.} The student-side trait gap ranks the transfer (Table~\ref{tab:lls}; Spearman
$+0.89$, $p=0.02$, six students). The teacher as its own student reaches $24\%$ dog mentions, close to the
$29.9\%$ the authors report for the same pair. Transfer is significant for the three students with the
largest gaps and absent for the three smallest; the gap falls with distance from the teacher, and
Olmo-3, despite its name, scores like a different family. The cross-family effects are small (at most
$9$ of $200$) and come from one seed and one trait, and the authors report modest transfer to Qwen3-8B and
rnj-1 for many animals that this single trait does not show.

\paragraph{What the elicit read sees here.} \Mtwo{} does not apply to this data: both sets are
human-written answers to their own questions. Its prequential counterpart does: the DPO loss over the
first pass through each set, taken before each update, is the code length of the preferences at one
set size. Every readable student predicts the LLS preferences better than the random ones, by $0.14$ to
$0.38$ nats per pair (teacher $0.316$ vs $0.691$; rnj-1 $0.558$ vs $0.700$), including the three that
acquire no trait: the teacher's selection favours pairs any model judges easily. Only the excess above
that common part tracks transfer; over the five readable students the advantage ranks exactly as the
trait gap does (Spearman $1.0$; Qwen2.5-7B's first pass was unstable and is not readable). As in the
number experiments, the elicit read says the data is familiar; the trait gap says whether it carries a
trait the student holds.

\section{Emergent misalignment}
\label{sec:em}

Emergent misalignment (EM) is the same pattern as SL: training on narrow data $A$ makes a broad
behaviour $B$ appear on unrelated questions \citep{betley2025em}. We measured EM on Qwen2.5-14B-Instruct
for $30$ published datasets with the recipe of \citet{turner2025organisms}; $8$ cause EM above $5\%$,
and the ranking replicates across three seeds (Spearman $0.98$--$0.99$). Scores of the training data
alone do not separate them, but scores of $B$ do.

\paragraph{The misaligned persona is held.} On the misaligned answers our fine-tunes produce, \Mtwo{}
reads $r=0.89$--$0.93$ of the model's own-answer gap on Qwen2.5-14B-Instruct, and $0.56$--$1.05$ on all
six models we tested (Qwen2.5 0.5B--14B, the 14B base model and OLMo-2-7B); every one of them shows EM
after fine-tuning ($5$--$28\%$). EM therefore elicits a persona the model already represents.

\paragraph{Which data elicits it is predictable.} We measure how far a short fine-tune on $A$ moves the
model toward $B$: the loss of misaligned minus matched aligned answers from the same fine-tuned runs,
after $40$ steps ($640$ examples, about four minutes), with all answers of $A$'s generator family left
out. Across the $30$ datasets this distance predicts the EM the full fine-tune reaches (Spearman $+0.57$,
AUC $0.78$ for EM above $5\%$), and $+0.55$ after partialling out the untouched model's loss on $A$; within
each generator family the order is nearly exact (Spearman $+0.62$ to $+1.00$). The same direction read
from the first update alone, before any training, gives $+0.64$ (AUC $0.76$; $+0.34$ beyond loss). Per
question, \Mtwo{} of misaligned answers on the untouched model tracks which questions misbehave most
(Spearman $+0.34$ within models over the $8$ standard questions). In elicit-versus-teach terms, \Mtwo{}
says the persona is held; narrow versus emergent depends on whether fitting $A$ recruits it. Training on
answers the model believes are false increases EM \citep{villegassuarez2026knowingly}; on our model the
on-policy version of their math experiment reproduces the direction weakly ($0\%$, $0\%$, $5\%$ for
correct, subtly wrong and absurd answers), and both our reads order the three sets correctly.

\bibliographystyle{plainnat}
\bibliography{refs}
\end{document}
